% Part: turing-machines
% Chapter: undecidability
% Section: enumerating-tms

\documentclass[../../../include/open-logic-section]{subfiles}

\begin{document}

\olfileid[pt]{tur}{und}{enu}
\olsection{Enumerando Máquinas de Turing}

\begin{explain}
Podemos mostrar que o conjunto de todas as máquinas de Turing é
!!{enumerable}. Isso decorre do fato de que cada máquina de Turing
pode ser finitamente descrita. O conjunto de estados e o vocabulário
da fita são conjuntos finitos. A função de transição é uma função
parcial de $Q \times \Sigma$ a $Q \times \Sigma \times \{\TMleft,
\TMright, \TMstay\}$ e, portanto, também pode ser especificada
listando-se seus valores para os finitos pares de argumentos para os
quais está definida.

Isso é verdade até certo ponto, mas há uma diferença sutil. A
definição de máquinas de Turing não fez restrição alguma sobre quais
!!{element}s o conjunto de estados e o alfabeto da fita podem ter.
Assim, por exemplo, para todo número real, há tecnicamente uma máquina
de Turing que usa esse número como um estado. Contudo, o
\emph{comportamento} da máquina de Turing é independente de quais
objetos servem como estados e vocabulário. Considere as duas máquinas
de Turing na \olref{fig:variants}.
\begin{figure}
\begin{center}
\begin{tikzpicture}[->,>=stealth',shorten >=1pt,auto,node distance=2.8cm,
                    semithick]
  \tikzstyle{every state}=[fill=none,draw=black,text=black]

  \node[initial,state]         (A)              {$q_0$};
  \node[state]         (B) [right of=A] {$q_1$};

  \path (A) edge [bend left] node {\TMtrans{\TMstroke}{\TMstroke}{\TMright}} (B)
        (B) edge [loop above] node {\TMtrans{\TMblank}{\TMblank}{\TMright}} (B)
            edge [bend left] node {\TMtrans{\TMstroke}{\TMstroke}{\TMright}} (A);
\end{tikzpicture}\\
\begin{tikzpicture}[->,>=stealth',shorten >=1pt,auto,node distance=2.8cm,
  semithick]
\tikzstyle{every state}=[fill=none,draw=black,text=black]

\node[initial,state]         (A)              {$s$};
\node[state]         (B) [right of=A] {$h$};

\path (A) edge [bend left] node {\TMtrans{A}{A}{\TMright}} (B)
(B) edge [loop above] node {\TMtrans{\TMblank}{\TMblank}{\TMright}} (B)
edge [bend left] node {\TMtrans{A}{A}{\TMright}} (A);
\end{tikzpicture}
\end{center}
\caption{Variantes da máquina \emph{Par}}
\ollabel{fig:variants}
\end{figure}
Esses dois diagramas correspondem a duas máquinas, $M$ com o alfabeto
de fita $\Sigma = \{\TMendtape,\TMblank,\TMstroke\}$ e conjunto de
estados $\{q_0,q_1\}$, e $M'$ com alfabeto $\Sigma' =
\{\TMendtape,\TMblank,A\}$ e estados $\{s,h\}$. Mas, quanto ao resto,
suas instruções são as mesmas: $M$ para em uma sequência de
$n$ $\TMstroke$ se, e somente se, $n$ é par, e $M'$ para em uma
sequência de $n$ $A$ se, e somente se, $n$ é par. Tudo o que fizemos
foi renomear $\TMstroke$ para~$A$, $q_0$ para~$s$ e $q_1$ para~$h$.
Esse exemplo se generaliza: podemos considerar máquinas de Turing como
sendo a mesma desde que uma resulte da outra por tal renomeação de
símbolos e estados. Na verdade, podemos simplesmente pensar nos
símbolos e estados de uma máquina de Turing como inteiros positivos:
em vez de $\sigma_0$ pense~$1$, em vez de $\sigma_1$ pense~$2$, etc.;
$\TMendtape$ é~$1$, $\TMblank$ é~$2$, etc. Dessa forma, a máquina
\emph{Par} torna-se a máquina representada na
\olref{fig:standard-even}.
\begin{figure}
\[\begin{tikzpicture}[->,>=stealth',shorten >=1pt,auto,node distance=2.8cm,
  semithick]
\tikzstyle{every state}=[fill=none,draw=black,text=black]

\node[initial,state]         (A)              {$1$};
\node[state]         (B) [right of=A] {$2$};

\path (A) edge [bend left] node {\TMtrans{3}{3}{\TMright}} (B)
(B) edge [loop above] node {\TMtrans{2}{2}{\TMright}} (B)
edge [bend left] node {\TMtrans{3}{3}{\TMright}} (A);
\end{tikzpicture}
\]
\caption{Uma máquina \emph{Par} padrão}
\ollabel{fig:standard-even}
\end{figure}
Poderíamos chamar de máquina \emph{padrão} uma máquina de Turing cujos
estados e símbolos são inteiros positivos, e considerar apenas
máquinas padrão de agora em diante.\footnote{A terminologia ``máquina
padrão'' não é padrão.}

Queríamos mostrar que o conjunto de máquinas de Turing é
!!{enumerable} e, tendo em mente as considerações acima, basta mostrar
que o conjunto de máquinas de Turing padrão é !!{enumerable}. Suponha
que nos seja dada uma máquina de Turing padrão $M = \tuple{Q, \Sigma,
q_0, \delta}$. Como poderíamos descrevê-la usando uma cadeia finita de
inteiros positivos? Primeiro listaremos o número de estados, os
próprios estados, o número de símbolos, os próprios símbolos e o
estado inicial. (Lembre-se, todos esses são inteiros positivos, já que
$M$ é uma máquina padrão.) E quanto a~$\delta$? O conjunto de
argumentos possíveis, isto é, pares $\tuple{q,\sigma}$, é finito, já
que $Q$ e~$\Sigma$ são finitos. Assim, a informação em~$\delta$ é
simplesmente a lista finita de todas as $5$-uplas $\tuple{q, \sigma,
q', \sigma', d}$ em que $\delta(q,\sigma) = \tuple{q', \sigma', D}$, e
$d$ é um número que codifica a direção~$D$ (digamos, $1$ para~$\TMleft$,
$2$ para~$\TMright$ e $3$ para~$\TMstay$).

Dessa forma, toda máquina de Turing padrão pode ser descrita por uma
lista finita de inteiros positivos, isto é, como uma sequência $s_M \in
(\PosInt)^*$. Por exemplo, a máquina \emph{Par} padrão é codificada
pela sequência
\[
2, \underbrace{1, 2}_Q, 3, \overbrace{1, 2, 3}^\Sigma, 1, \underbrace{1, 3, 2, 3, 2}_{\delta(1,3) = \tuple{2,3,R}}, 
\overbrace{2, 2, 2, 2, 2}^{\delta(2,2) = \tuple{2,2,R}},
\underbrace{2, 3, 1, 3, 2}_{\delta(2,3) = \tuple{1,3,R}}.
\]
\end{explain}

\begin{thm}
Há funções de $\Nat$ a~$\Nat$ que não são Turing-computáveis.
\end{thm}

\begin{proof}
Sabemos que o conjunto de sequências finitas de inteiros
positivos~$(\PosInt)^*$ é !!{enumerable}
(\cref{sfr:siz:zigzag:prob:posint-star}). Isso nos dá que o conjunto
de descrições de máquinas de Turing padrão, como um subconjunto
de~$(\PosInt)^*$, é ele próprio enumerável. Toda função
Turing-computável de $\Nat$ a~$\Nat$ é computada por alguma (na
verdade, muitas) máquina de Turing. Renomeando seus estados e símbolos
para inteiros positivos (em particular, $\TMendtape$ como~$1$,
$\TMblank$ como~$2$ e $\TMstroke$ como~$3$), podemos ver que toda
função Turing-computável é computada por uma máquina de Turing padrão.
Isso significa que o conjunto de todas as funções Turing-computáveis de
$\Nat$ a~$\Nat$ é também enumerável.

Por outro lado, o conjunto de todas as funções de $\Nat$ a~$\Nat$ não
é !!{enumerable} (\cref{sfr:siz:red:prob:nat-nat}). Se todas as funções
fossem computáveis por alguma máquina de Turing, poderíamos enumerar o
conjunto de todas as funções listando todas as descrições de máquinas
de Turing que as computam. Logo, há algumas funções que não são
Turing-computáveis. 
\end{proof}

\begin{prob}
  Você consegue pensar em uma maneira de descrever máquinas de Turing
  que não exija que os estados e os símbolos do alfabeto sejam
  listados explicitamente? Você pode definir sua própria noção de
  máquina ``padrão'', mas diga algo sobre por que toda máquina de
  Turing pode ser computada por uma máquina ``padrão'' no seu novo
  sentido.
\end{prob}

\end{document}
